\section{Benchmark Design}\label{sec-4}

\subsection{Five Capability Pillars}\label{sec-4-1}

We decompose rare-disease diagnosis into five distinct capability pillars, each individually testable, jointly forming the evaluation surface. The pillar choice is grounded in three considerations: (i) clinical workflow steps (phenotype recognition \ensuremath{\rightarrow} DDx generation \ensuremath{\rightarrow} genetic confirmation \ensuremath{\rightarrow} family interpretation \ensuremath{\rightarrow} communication), (ii) the dimensions on which existing rare-disease agents claim differentiation (DeepRare's reference accuracy, MAI-DxO's budgeted reasoning), and (iii) reviewer-anticipated criticism that accuracy alone is insufficient \citep{equitymedqa2024}; Pillar 5 directly addresses this.

The five pillars, and which data layer carries each, are summarised in the
benchmark-surface figure. Briefly: P1 Phenotype Extraction (free-text EHR
\ensuremath{\rightarrow} ranked HPO list; P/R/F1); P2 Phenotype-only DDx (HPO list \ensuremath{\rightarrow} ranked
diseases; Recall@1/3/5/10, median rank, MRR); P3 Genotype-aware DDx (HPO +
structured variants \ensuremath{\rightarrow} disease + causative gene; Gene Top-k, cross-mapped R@1);
\textbf{P4} Family-aware DDx (HPO + trio/pedigree \ensuremath{\rightarrow} disease + mode of inheritance;
deferred to v2 pending pedigree-bearing data); and P5 Reasoning
Faithfulness (agent trace + prediction \ensuremath{\rightarrow} a four-axis factual / relevance /
depth / faithfulness score; the current 40-trace LLM-judge sensitivity study
is exploratory, and physician validation remains future work).

\textbf{Why five and not three.} P1 and P5 in particular face the objection ``isn't this just preprocessing or just a stylistic axis?''. Both are load-bearing for our central claim. P1 is independent because \Cref{sec-7-1} shows, in a same-case paired test, that downstream P2 R@1 is measurably sensitive to HPO-extraction quality --- extraction is not a free preprocessing step (and the dual-pass design is what lets this be measured without confounding it with dataset difficulty). P5 is independent because \Cref{sec-7-4}/\Cref{sec-7-5} examine whether faithfulness ranking decouples from accuracy ranking (repository-plan Hypothesis 10); we report H10 as exploratory (a modest same-trace judge difference remains, but the current Gemini-to-Claude comparison cannot identify a pure family effect), and the durable point is that accuracy-only evaluation can miss hallucinated citations and unfaithful reasoning chains --- a Type-1 error category discussed by DeepRare \citep{deeprare2026}.

\textbf{Bias as cross-cutting lens, not pillar.} Following HELM and MedHELM
precedent \citep{helm2022,medhelm2025}, we apply available stratifications
(prevalence tier, specialty, age group, sex, and HPO density) across relevant
pillar metrics rather than treating bias as a separate axis. Dedicated
fairness evaluations \citep{equitymedqa2024,heal2024,omiye2023,zack2024}
remain complementary: ancestry and language analyses cannot be claimed in v1
because the current frozen resources do not support them.

\subsection{Datasets --- Three Diagnostic Layers + One Structured-EHR Probe}\label{sec-4-2}

We assemble three diagnostic layers and one secondary structured-EHR probe; their
sources and sizes appear as the columns of the benchmark-surface figure, and
the full per-layer breakdown (disease counts, ID anchors, free-text / gold-HPO
/ variant availability) is detailed in Appendix A. The diagnostic layers are L1
Phenotype Backbone (Phenopacket-Store + RareBench HF; 11,173 cases; HPO-only,
gold HPO, variants on PP-Store), L2 Scale + Free Text (RareArena RDS/RDC;
72,661 verbatim case reports), and L4 Cutoff-After Holdout (self-built PMC
OA, publication date \ensuremath{\geq} 2024-01-01; 198 model-verified cases with physician
annotation in progress). Separately,
\textbf{S-EHR} is a credentialed MIMIC-IV-3.1 probe \citep{mimiciv31}
(956 admissions, 239 exact-mapped
ORPHA labels) for early structured-event prediction and ICD leakage auditing.
The replacement protocol is specified but not yet scored; it is not pooled
with differential-diagnosis results.

\textbf{Rationale per layer.}

\begin{itemize}
\tightlist
\item
  \textbf{L1} establishes apples-to-apples comparison with RareBench's KDD'24 numbers --- required by reviewers (``why not just compare to RareBench'').
\item
  \textbf{S-EHR} is designed to test whether agents can use real structured hospital events rather than curated case reports. The 24-hour primary snapshot uses timestamped labs, medications, procedures and services; it excludes target-bearing diagnosis codes/titles, post-window events and free text. Gold is derived from exact ICD-10\ensuremath{\rightarrow}Orphanet mapping, so this is code-supervised retrospective prediction, not independently adjudicated diagnosis. A paired audit will compare ICD titles, ICD codes only, and context after removing target-bearing entries.
\item
  \textbf{L2} addresses scale; RareArena's 72,661 cases span 45.6\% of Orphanet. Released CC-BY-NC-SA, so we use it for academic evaluation only and acknowledge license bounds.
\item
  \textbf{L4} responds to the leakage and benchmark-composition concerns in the
  2026 systematic review \citep{sysreview2026}. We extract PMC OA case reports
  published after 2024-01-01 via E-utilities and map diagnoses to Orphanet.
  The 198-case post-cutoff set has model-based verification but physician
  annotation is still in progress. Exact PMCIDs overlap RareArena, so L4 is a
  temporal sensitivity analysis, not an independently clean holdout.
\end{itemize}

\textbf{Total resource pool: \textasciitilde85,000 cases}, with v1 analyses stratified where
supported by prevalence tier, specialty, age group, sex, and HPO density. The
current scored diagnostic study is English-language; a Chinese diagnostic
layer is deferred (\Cref{sec-9}).

\textbf{Evaluation N per dataset --- honest disclosure.} We deliberately separate
``pool size'' (the released benchmark) from ``evaluation N'' (what we run for the
v1 paper).
- \textbf{RareBench-HF} has 1,122 cases. Primary general-agent cells use the full
attempted N; DeepRare and MAI-DxO have explicitly smaller adapter-specific N.
- \textbf{Large layers} (Phenopacket-Store 10,051; RareArena RDS 72,661): the primary
\texttt{llm\_control}/MDAgents/MedAgents matrix uses a shared, prevalence-stratified
cap of N=2,000 case IDs (seed=42). DeepRare, MAI-DxO, and a small number of
failed-to-launch adapter cells have smaller attempted N shown in \Cref{tbl:main}.
- \textbf{MIMIC} will use its own attempted-admission denominator and separate
receipts after the replacement experiment is run.

This is a prevalence-stratified evaluation, not a power-stratified extrapolation to full pool; bootstrap CIs in \Cref{sec-6} quantify the resulting uncertainty per cell.
We confirmed via our sampling-validation check that the N=2,000 stratified sample reproduces the full-N
prevalence band and HPO-organ-system distribution within \ensuremath{\pm}2 pp.

\subsection{Canonical Case Representation}\label{sec-4-3}

The three diagnostic layers ingest into a single Pydantic v2 schema,
\texttt{CanonicalCase} (illustrated in the canonical-case figure), and every agent
adapter projects from this representation to the agent's native input. The
pending MIMIC protocol uses an analogous model-input/evaluation-only record and
maps the structured snapshot at the adapter boundary. The record groups an
identity block (\texttt{case\_id}, \texttt{source\_dataset}, \texttt{source\_split}, \texttt{language}), an
input block (demographics; \texttt{free\_text\_vignette} / \texttt{synthetic\_vignette};
\texttt{gold\_hpo\_terms}; \texttt{variants}; local \texttt{vcf\_path}; \texttt{family}), a parallel-ID gold
label, and free-form metadata.

Three design decisions deserve note: (a) gold labels are parallel IDs (OMIM/ORPHA/CCRD), reflecting genuine ontology disagreement across datasets --- Phenopacket-Store uses OMIM, RareArena uses ORPHA, CCRD anchors the Chinese-listed 207 diseases. Evaluator must accept cross-mapped matches via Orphadata (\Cref{sec-4-5}). (b) \texttt{synthetic\_vignette} is distinguished from \texttt{free\_text\_vignette} so we can audit any evaluation that relies on LLM-synthesized prose (\Cref{sec-7-2} disclosure). (c) \texttt{vcf\_path} carries a local-only file pointer; PhysioNet DUA prohibits transmission of identifiable EHR data to external LLM APIs, so adapter shims projecting Pillar 3 inputs convert structured variant info to abstracted strings before any cloud-LLM call.

\subsection{Dual-pass evaluation}\label{sec-4-4}

Every pillar is evaluated in two modes. In Pass A (gold-HPO, primary) the
case's curated HPO terms are fed directly to the agent's downstream pillar,
bypassing its own extractor --- isolating downstream capability and enabling
apples-to-apples comparison on identical inputs. In Pass B (end-to-end) the
raw free-text vignette is fed in and the agent extracts HPO itself before
reasoning, measuring real deployment performance. The Pass A \ensuremath{-} Pass B delta
is itself a reported metric (following RareBench's phenotype-vs-EHR-text
comparison \citep{rarebench2024}); we show it is non-uniform across agents
(\Cref{sec-7-1}), which turns the input heterogeneity of a mixed-agent lineup into a
measured axis rather than a confound (\Cref{sec-5-1}).

\subsection{Metrics and matching}\label{sec-4-5}

We report Recall@1/3/5/10, median rank and MRR for the DDx pillars, P/R/F1 for
phenotype extraction, and task-success rate; a recommended tier (pass\^{}k,
cost-normalised accuracy, calibration, reference accuracy) and an exploratory
tier (step-level and reasoning-faithfulness scores) are defined in Appendix C.
Every accuracy metric is stratified where metadata and sample size permit;
v1 reports prevalence tier, specialty, age group, sex, and HPO-density
analyses, while ancestry and language remain deferred. A predicted
disease ID counts as a hit if it is prefix-equal on the same ontology,
cross-mapped via Orphadata (OMIM \ensuremath{\leftrightarrow} ORPHA), or fuzzy-matches an Orphanet name or
synonym at rapidfuzz score \ensuremath{\geq} 90 --- a threshold calibrated against 217 audited
borderline cases (\textgreater85\% of the 70--89 band were false positives); the audit tape
is released for reproducibility (Appendix N).


\begin{figure*}[t]
\centering
\includegraphics[width=0.82\textwidth]{figures/fig1_overview.png}
\caption{\textbf{RareAgentBench overview.} The three diagnostic layers use a shared \texttt{CanonicalCase} contract; the separately reported MIMIC replacement probe preserves the same model-input/gold separation and is projected at the adapter boundary. Eight published systems plus one no-scaffold control run through subprocess-isolated shims. Diagnostic pillars use the two-pass protocol; the pending MIMIC experiment is specified to use timestamped structured events and an ICD leakage audit.}\label{fig:fig1_overview}
\end{figure*}

